% TOP_PROBLEM: {"id": 1452, "title": "Farrell-Jones Conjecture", "category_id": 4, "category": {"id": 4, "name": "algebra", "display_name": "Algebra", "description": "Group theory, ring theory, field theory, and algebraic structures.", "slug": "algebra", "order_index": 4, "created_at": "2026-07-31T15:26:25.671Z"}, "status": "open"}
% FIRST_POSTED: 2026-09-27
% !TeX program = lualatex
% Compile twice: lualatex 1452.tex
% Original entry retained; research subsection and [E] references added.
% No external bibliography, image, data, or style file is required.
\documentclass[11pt,a4paper]{article}
\usepackage[margin=24mm,headheight=14pt]{geometry}
\usepackage{fontspec}
\setmainfont{Latin Modern Roman}
\setsansfont{Latin Modern Sans}
\setmonofont{Latin Modern Mono}
\usepackage{amsmath,amssymb,mathtools,amsthm}
\usepackage{unicode-math}
\setmathfont{Latin Modern Math}
\usepackage{microtype}
\usepackage{enumitem,xurl,needspace}
\usepackage[dvipsnames]{xcolor}
\usepackage{hyperref}
\hypersetup{unicode=true,colorlinks=true,linkcolor=MidnightBlue,
 urlcolor=MidnightBlue,bookmarksnumbered=true,
 pdftitle={Research notebook - Problem 1452},
 pdfauthor={Research attempt; not independently refereed}}
\usepackage{bookmark}
\usepackage{titlesec}
\titleformat{\section}{\Large\bfseries\raggedright}{\thesection}{0.7em}{}
\usepackage{fancyhdr}
\pagestyle{fancy}\fancyhf{}
\fancyhead[L]{\small\sffamily Research notebook}
\fancyhead[R]{\small\sffamily Problem 1452}
\fancyfoot[L]{\scriptsize 27 September 2026}
\fancyfoot[C]{\thepage}
\fancyfoot[R]{\scriptsize Unrefereed research attempt}
\setlength{\parindent}{0pt}
\setlength{\parskip}{4pt plus 1pt}
\setlength{\emergencystretch}{3em}
\clubpenalty=10000
\widowpenalty=10000
\displaywidowpenalty=10000
\setcounter{secnumdepth}{1}
\setlist[itemize]{leftmargin=3em,itemsep=2pt,topsep=3pt,parsep=0pt,before=\small}
\allowdisplaybreaks
\newcommand{\N}{\mathbb{N}}
\newcommand{\Z}{\mathbb{Z}}
\newcommand{\Q}{\mathbb{Q}}
\newcommand{\R}{\mathbb{R}}
\newcommand{\C}{\mathbb{C}}
\newcommand{\F}{\mathbb{F}}
\newcommand{\A}{\mathbb{A}}
\newcommand{\PP}{\mathbb P}
\newcommand{\E}{\mathbb{E}}
\newcommand{\eps}{\varepsilon}
\newcommand{\dd}{\,\mathrm d}
\newcommand{\sref}[2]{\hyperlink{source:#1:#2}{[S#2]}}
\newcommand{\eref}[2]{\hyperlink{extra:#1:#2}{[E#2]}}
\newif\ifshowcatalogue
\showcataloguetrue
\newcommand{\cataloguescope}[1]{\ifshowcatalogue
 \paragraph{Exact scope in the supplied catalogue.}
 \begin{quote}\small #1\end{quote}\fi}
\newtheorem{notebooklemma}{Lemma}[section]
\newtheorem{notebookproposition}[notebooklemma]{Proposition}
\numberwithin{equation}{section}
\begin{document}
\setcounter{section}{0}
% ================================================================
% problemId: 1452
% Canonical problem ID: 1452
% exactTarget SHA-256: ea425e3c220d6073d10c5075820287000d53c720087ccf8b31cd8e38fbd20f94
\clearpage
\section*{\texorpdfstring{Farrell-Jones Conjecture}{Farrell-Jones Conjecture}}\label{1452}
\subsection{Definitions and mathematical statement}
A group is virtually cyclic if it contains a cyclic subgroup of finite index. Let $E_{\rm VCyc}G$ be the classifying $G$-space for that family, characterized by its subgroup fixed-point sets. Using the algebraic $K$- and $L$-theory spectra with the source's standard decorations, the required maps are
\[
 H_n^G(E_{\rm VCyc}G;\mathbf K_{\Z})\longrightarrow K_n(\Z G),\qquad
 H_n^G(E_{\rm VCyc}G;\mathbf L_{\Z})\longrightarrow L_n(\Z G).
\]
They are induced by $E_{\rm VCyc}G\to\mathrm{pt}$ and are conjectured to be isomorphisms for every group and degree. Coefficients other than the integral group ring are not added to the selected statement.
\par\smallskip\textit{Formulation and concept references:} \sref{1452}{1}.
\cataloguescope{Prove the Farrell-Jones conjecture for every group G: the assembly maps from the equivariant homology of the classifying space for the family of virtually cyclic subgroups to the algebraic K-theory, respectively L-theory, of the group ring \ensuremath{\mathbb Z}G are isomorphisms.}
\subsection{Short English statement}
Can the algebraic $K$- and $L$-theory of every integral group ring be assembled completely from its virtually cyclic subgroups?
% BEGIN RESEARCH INSERTION
\subsection{Research attempt}

\paragraph{Editorial note on source scope.}
The original statement is retained. Lück's Conjectures 8.7--8.8 allow more general coefficient categories than the selected integral-group-ring target. We pursue their specialization to $\Z G$, with ultimate decoration $\langle-\infty\rangle$ on the $L$-theory side, not a silently strengthened coefficient conjecture \sref{1452}{1}, Section 8.4.

\paragraph{Current frontier.}
Hyperbolic and finite-dimensional CAT(0) groups are among the established classes, with extensive inheritance properties \sref{1452}{1}, Theorem 8.12. Andrew--Guerch--Hughes obtain results about automorphisms of relatively hyperbolic groups with specified peripheral hypotheses, not arbitrary groups \eref{1452}{2}. For the argument below we use continuity under directed unions of subgroups, established for integral algebraic $K$-theory and ultimate $L$-theory by Bartels--Echterhoff--Lück \eref{1452}{1}, Sections 2--3 and Lemma 5.2. An inverse limit of finite quotients is a different operation.

\paragraph{Attack route.}
Controlled topology seeks to turn a finite-support class into one assembled from virtually cyclic subgroups. A counterexample construction seeks a class whose defect persists at every later stage. We make this persistence requirement exact. We then test whether finite quotients or completed group rings could supply the missing repair, and find that compatible approximate inverses already fail to produce an inverse in an elementary group ring.

\paragraph{Attempt: persistence of a finite-stage defect.}
Fix a degree $n\in\Z$ and one of the two theories. Write
\[
 A_n(G)=H_n^G(E_{\rm VCyc}G;\mathbf T_{\Z}),\qquad
 B_n(G)=T_n(\Z G),\qquad \mu_G:A_n(G)\to B_n(G),
\]
where $T$ is $K$ or $L^{\langle-\infty\rangle}$. Let $G=\bigcup_{i\in I}G_i$ be a directed union of subgroups. The needed continuity statements are
\begin{equation}\label{r68:colimit}
 \varinjlim_i A_n(G_i)\simeq A_n(G),\qquad
 \varinjlim_i B_n(G_i)\simeq B_n(G),
\end{equation}
compatibly with assembly. On the source side, restricting $E_{\rm VCyc}G$ to $G_i$ gives a model for $E_{\rm VCyc}G_i$: virtual cyclicity is an intrinsic property of the subgroup. On the target side, the finite algebraic data defining classes and relations give continuity; the all-degree and $L$-decoration assertions are supplied by \eref{1452}{1}, Lemma 5.2, rather than presumed from a $K_0$ analogy.

\begin{notebooklemma}\label{r68:persist}
Under~\eqref{r68:colimit}, a failure of surjectivity of $\mu_G$ is represented by a class $y_i\in B_n(G_i)$ whose image is outside $\operatorname{im}\mu_{G_j}$ for every $j\ge i$. A failure of injectivity is represented, after possibly moving to a later stage, by nonzero classes $x_j\in\ker\mu_{G_j}$ at every subsequent stage.
\end{notebooklemma}
\begin{proof}
Choose $y\in B_n(G)\setminus\operatorname{im}\mu_G$ and represent it by $y_i$. If at any stage $j\ge i$ its image were $\mu_{G_j}(a_j)$, then the image of $a_j$ in the colimit would lift $y$, a contradiction.

For injectivity, choose $0\ne x\in A_n(G)$ with $\mu_G(x)=0$ and represent it by $x_i$. An element that becomes zero in a filtered colimit of abelian groups becomes zero at some later stage. Apply this to $\mu_{G_i}(x_i)$ and move to that stage. At every subsequent stage its assembly image remains zero, while its source image cannot vanish because its image $x$ in the colimit is nonzero.
\end{proof}

Taking all finitely generated subgroups as the directed family shows that any counterexample would already force a counterexample among finitely generated groups. Conversely, proving the target for every finitely generated group suffices for all groups. This does not claim reduction to finite groups or to finitely presented groups; finite generation does not imply finite presentation.

\paragraph{An exact eventual-repair criterion.}
For a proposed directed presentation of $G$, the missing conditions can be written without assuming that every stage satisfies the conjecture:
\begin{align}
 &\forall i\ \forall y_i\in B_n(G_i)\ \exists j\ge i\ \exists a_j\in A_n(G_j):
       \mu_{G_j}(a_j)=(y_i)_j; \label{r68:repair1}\\
 &\forall i\ \forall x_i\in\ker\mu_{G_i}\ \exists j\ge i:
       (x_i)_j=0. \label{r68:repair2}
\end{align}
These conditions are equivalent to $\mu_G$ being an isomorphism. Sufficiency follows from the same finite-stage argument: an eventual kernel relation is first moved to a stage where it is actually a kernel relation, and then~\eqref{r68:repair2} applies. Conversely, if $\mu_G$ is an isomorphism, a lift in the colimit is represented at a finite stage, and equality of its assembly image with a prescribed $y_i$ holds at some later common stage. The injective case follows by eventual vanishing.

This is an equivalence obtained from continuity, not a new proof of either repair condition. Its research use would be a uniform controlled construction producing $j$ and $a_j$, or a class that prevents all such choices. No such construction is supplied for arbitrary finitely generated groups.

\paragraph{Attempted repair by finite quotients: a vacuous test.}
Every finite group is virtually cyclic, since it contains the trivial cyclic subgroup with finite index. For such a group, the virtually cyclic classifying space may be a point and assembly is an identity. Consequently checking the target on all finite quotients of a residually finite group supplies no additional restriction at all: those checks automatically succeed. Residual finiteness is an inverse approximation property, not the directed-union hypothesis in~\eqref{r68:colimit}.

Even compatible finite-order algebraic solutions need not reconstruct an algebraic solution in the original ring. In
\[
 R=\Z[t,t^{-1}]=\Z[\Z],\qquad s=t-1,
\]
consider $u=2-t=1-s$. For every $r\ge1$,
\begin{equation}\label{r68:geometric}
 (1-s)\sum_{j=0}^{r-1}s^j=1-s^r.
\end{equation}
Thus $u$ has compatible inverses modulo $(t-1)^r$, and its formal inverse exists in the $(t-1)$-adic completion. Nevertheless $u$ is not a unit in $R$. To prove this, define the width of a nonzero Laurent polynomial to be its highest exponent minus its lowest. Over the domain $\Z$, widths add under multiplication. A product equal to one forces both widths to be zero; its factors must be monomials, with coefficients $\pm1$. The two-term polynomial $2-t$ is not of that form.

The quotient rings here are nilpotent-order approximations, not claimed to be integral group rings of finite quotient groups. The example disproves the algebraic inference from compatible inverses in approximations to an inverse in the original group ring. It is not a Farrell--Jones counterexample: $\Z$ itself is virtually cyclic, so its target assembly is already an identity.

\paragraph{Stress test.}
A class with no lift at an early stage may acquire one later; persistence, not a single unsuccessful computation, is needed for a global counterexample. A null relation can also appear only after adjoining further group elements. Both issues are handled by the filtered-colimit quantifiers. Conversely, substituting an inverse limit or completion is unjustified without an additional reconstruction theorem. The explicit compatible inverse in~\eqref{r68:geometric} shows why such a theorem cannot be inferred from compatibility alone.

The script verifies the geometric identity for $1\le r\le30$ by exact symbolic arithmetic. The identity and the nonunit argument hold for all $r$ by the displayed algebraic proofs. No finite computation of assembly groups is claimed. Coefficient-category inheritance results in the sources are used only where they specialize to the fixed integral coefficients; they are not added to the conclusion.

\paragraph{Outcome.}
\textbf{Outcome: Exploratory approach with unresolved gap.}

The notebook makes finite-stage persistence and eventual repair precise, and proves that the proposed approximation shortcut cannot be justified from compatible algebraic solutions alone. These are deductions and obstructions around established continuity, not a new class of groups satisfying Farrell--Jones. A general controlled repair theorem or a genuinely persistent finitely generated counterexample is still required.

% END RESEARCH INSERTION
\subsection{Sources}
\begin{itemize}\raggedright
\item[\textnormal{[S1]}]\hypertarget{source:1452:1}{} Wolfgang Lück, Survey on the Farrell-Jones Conjecture, Conjectures 8.7 and 8.8, 2025.
\url{https://arxiv.org/html/2507.11337}.
{\small\textit{Catalogue formulation source.}}
\item[\textnormal{[S2]}]\hypertarget{source:1452:2}{} Automorphisms of relatively hyperbolic groups and the Farrell-Jones conjecture.
\url{https://link.springer.com/article/10.1007/s00208-026-03431-7}.
{\small\textit{Catalogue research/status reference; not a proof endorsement.}}
\item[\textnormal{[S3]}]\hypertarget{source:1452:3}{} On the Farrell-Jones conjecture for localising invariants.
\url{https://arxiv.org/abs/2111.02490}.
{\small\textit{Catalogue research/status reference; not a proof endorsement.}}

% BEGIN ADDED RESEARCH REFERENCES
\item[\textnormal{[E1]}]\hypertarget{extra:1452:1}{} Arthur Bartels, Siegfried Echterhoff, and Wolfgang Lück, \emph{Inheritance of Isomorphism Conjectures under colimits}, arXiv:math/0702460 (2007), Sections 2--3 and Lemma 5.2.
\url{https://arxiv.org/html/math/0702460}.
{\small\textit{Continuity of assembly and of algebraic $K$- and ultimate $L$-theory; the notebook uses directed unions. Consulted 27 September 2026.}}
\item[\textnormal{[E2]}]\hypertarget{extra:1452:2}{} Naomi Andrew, Yassine Guerch, and Sam Hughes, \emph{Automorphisms of relatively hyperbolic groups and the Farrell--Jones conjecture}, Mathematische Annalen 395, Article 89 (2026), DOI 10.1007/s00208-026-03431-7.
\url{https://link.springer.com/article/10.1007/s00208-026-03431-7}.
{\small\textit{Specific relatively hyperbolic and peripheral hypotheses; frontier comparison, not a universal theorem. Consulted 27 September 2026.}}
% END ADDED RESEARCH REFERENCES
\end{itemize}

\end{document}
